% ============================================================
% ANEXO A — Detalle Financiero
% ============================================================
\chapter{Detalle Financiero}
\label{anx:financiero}

Este anexo recoge el desglose detallado de la inversi\'on realizada en
el Proyecto Demeter, complementando las tablas resumen del
cap\'itulo~\ref{cap:sostenibilidad}.

\section{Inventario de hardware adquirido}
\label{sec:inventario}

El inventario se presenta en dos bloques seg\'un su trazabilidad
documental. La \textbf{instalaci\'on hidr\'aulica y de potencia} procede
de un \'unico expediente de compra con presupuesto formal y se detalla
unidad a unidad. La \textbf{electr\'onica de control} se adquiri\'o de
forma incremental en distribuci\'on \textit{online} durante el
desarrollo, sin expediente unificado, por lo que se inventar\'ia por
tipo de componente y con precio de referencia de mercado.

\begin{table}[H]
    \centering
    \caption{Instalaci\'on hidr\'aulica y de potencia. Expediente
             \'unico con trazabilidad completa.}
    \label{tab:inventario_hidraulico}
    \begin{tabular}{p{5.2cm}crrp{2.6cm}}
        \toprule
        \textbf{Componente} & \textbf{Uds.} & \textbf{\euro/ud.} & \textbf{Total} & \textbf{Ubicaci\'on actual} \\
        \midrule
        Bomba el\'ectrica autocebada 12\,V (Coolty) & 4 & 29,99 & 119,96\,\euro & Almac\'en \mbox{ESIBot} \\
        Electrov\'alvula DC 12\,V 1\,\textquotedbl  & 6 & 36,74 & 220,44\,\euro & Almac\'en \mbox{ESIBot} \\
        V\'alvula reguladora de presi\'on 1\,\textquotedbl & 4 & 31,91 & 127,64\,\euro & Almac\'en \mbox{ESIBot} \\
        V\'alvula antirretorno 1\,\textquotedbl     & 4 & 18,11 & 72,44\,\euro  & Almac\'en \mbox{ESIBot} \\
        Barril 60\,L boca ancha HDPE sin BPA        & 4 & 28,45 & 113,80\,\euro & Almac\'en \mbox{ESIBot} \\
        Kit de grifos para tanque                   & 2 & 16,80 & 33,60\,\euro  & Almac\'en \mbox{ESIBot} \\
        Fuente carril DIN 12\,V 120\,W              & 1 & 29,93 & 29,93\,\euro  & Almac\'en \mbox{ESIBot} \\
        Cable plano gemelo 2 n\'ucleos 10\,m        & 1 & 25,99 & 25,99\,\euro  & Almac\'en \mbox{ESIBot} \\
        \midrule
        Base imponible                              & \multicolumn{2}{r}{}       & 743,80\,\euro & \\
        IVA (21\,\%)                                & \multicolumn{2}{r}{}       & 156,20\,\euro & \\
        \midrule
        \textbf{TOTAL EXPEDIENTE}                   & \multicolumn{2}{r}{}       & \textbf{900,00\,\euro} & \\
        \bottomrule
    \end{tabular}
\end{table}

\textbf{Trazabilidad del expediente}: presupuesto n\textsuperscript{o}
\textbf{26317}, emitido el \textbf{31/10/2025} por \textit{Juan Rodrigo
Jim\'enez-D\'iaz --- MR MICRO Sevilla} (NIF ES45659910B, C/ Feria 135,
Local UPI, 41002 Sevilla) a nombre de la \textbf{Escuela T\'ecnica
Superior de Ingenier\'ia Agron\'omica} (Carretera de Utrera km 1, 41013
Sevilla; NIF Q4118001I), bajo la referencia \textbf{ESIBOT}. Forma de
pago: transferencia bancaria. Material recibido y disponible desde
\textbf{enero de 2026}.

\begin{table}[H]
    \centering
    \caption{Electr\'onica de control y sensores.}
    \label{tab:inventario_electronico}
    \begin{tabular}{p{5.2cm}crrp{2.6cm}}
        \toprule
        \textbf{Componente} & \textbf{Uds.} & \textbf{\euro/ud.} & \textbf{Total} & \textbf{Ubicaci\'on actual} \\
        \midrule
        ESP32-S3-DevKitC-1                     & 3  & $\sim$15 & 45\,\euro & Almac\'en \mbox{ESIBot} \\
        Raspberry Pi 4B (4\,GB)                & 1  & $\sim$75 & 75\,\euro & Almac\'en \mbox{ESIBot} \\
        MicroSD 32\,GB clase 10                & 1  & $\sim$9  & 9\,\euro  & Almac\'en \mbox{ESIBot} \\
        Sensor DS18B20 (sonda impermeable)     & 10 & $\sim$4  & 40\,\euro & Almac\'en \mbox{ESIBot} \\
        Sensor capacitivo de humedad de suelo  & 10 & $\sim$3  & 30\,\euro & Almac\'en \mbox{ESIBot} \\
        Sensor DHT22                           & 10 & $\sim$5  & 50\,\euro & Almac\'en \mbox{ESIBot} \\
        M\'odulo rel\'e 4 canales              & 3  & $\sim$6  & 18\,\euro & Almac\'en \mbox{ESIBot} \\
        Convertidores TTL/USB                  & 2  & $\sim$5  & 10\,\euro & Almac\'en \mbox{ESIBot} \\
        Fuentes de alimentaci\'on 5\,V         & 3  & $\sim$5  & 15\,\euro & Almac\'en \mbox{ESIBot} \\
        Bater\'ias LiPo 3000\,mAh              & 2  & $\sim$12 & 24\,\euro & Almac\'en \mbox{ESIBot} \\
        Carcasas/protecciones IP65             & 2  & $\sim$10 & 20\,\euro & Almac\'en \mbox{ESIBot} \\
        Cables, jumpers, \textit{breadboard}   & \multicolumn{2}{r}{lote} & 30\,\euro & Almac\'en \mbox{ESIBot} \\
        \midrule
        \textbf{TOTAL ELECTR\'ONICA}           & \multicolumn{2}{r}{}     & \textbf{366\,\euro} & \\
        \bottomrule
    \end{tabular}
\end{table}

\textbf{Trazabilidad de la electr\'onica}: adquirida en distribuci\'on
\textit{online} en pedidos peque\~nos a lo largo de 2025 y 2026, sin
expediente de compra unificado ni n\'umero de factura \'unico. Los
importes son precios de referencia de mercado por unidad. El material
est\'a \'integramente depositado en el almac\'en de la asociaci\'on.

\textbf{Identificaci\'on de equipos}. El nodo \textit{gateway} es
identificable de forma un\'ivoca por su direcci\'on MAC
\texttt{9C:13:9E:A8:6F:CC}, que aparece codificada en el firmware de los
nodos sensores como destino ESP-NOW. Los nodos restantes se identifican
por su \texttt{NODE\_ID} de compilaci\'on (1 \textit{gateway},
2 sensor, 3 actuador).

\textbf{Inversi\'on total en hardware}: 900,00\,\euro\ (hidr\'aulica) +
366\,\euro\ (electr\'onica) = \textbf{1.266\,\euro}.

\section{Hist\'orico de gastos en infraestructura cloud}

\begin{table}[H]
    \centering
    \caption{Gastos mensuales en Google Cloud Platform.}
    \label{tab:cloud_historico}
    \begin{tabular}{lr}
        \toprule
        \textbf{Mes} & \textbf{Importe (\textbf{\euro})} \\
        \midrule
        01/2026 & 25\,\euro \\
        02/2026 & 25\,\euro \\
        03/2026 & 25\,\euro \\
        04/2026 & 25\,\euro \\
        05/2026 & 25\,\euro \\
        06/2026 & 25\,\euro \\
        07/2026 & 25\,\euro \\
        08/2026 & 25\,\euro \\
        \midrule
        \textbf{TOTAL ACUMULADO} & \textbf{200\,\euro} \\
        \bottomrule
    \end{tabular}
\end{table}

El gasto se mantuvo constante en 25\,\euro/mes porque la carga del
sistema no variaba: una \'unica m\'aquina virtual encendida de forma
permanente, con un volumen de telemetr\'ia muy por debajo de cualquier
umbral de facturaci\'on por uso. La cifra de 200\,\euro\ es
\textbf{definitiva}: la cuenta ya ha sido cerrada y no habr\'a m\'as
cargos.

Conviene subrayar el dato m\'as revelador de esta tabla: \textbf{el reloj
de facturaci\'on no se detuvo cuando lo hizo el desarrollo}. El \'ultimo
\textit{commit} del repositorio es de marzo de 2026, pero la
infraestructura sigui\'o facturando \textbf{cinco meses m\'as} sin ning\'un
trabajo asociado --- 125\,\euro, el 62\,\% del total gastado en cloud, se
consumieron despu\'es de que el proyecto dejara de avanzar. Es la
ilustraci\'on m\'as directa del argumento del
cap\'itulo~\ref{cap:sostenibilidad}: el coste de un sistema desplegado es
recurrente aunque el proyecto no lo sea, y solo deja de serlo cuando
alguien lo apaga de forma deliberada.

Las facturas mensuales originales y los datos identificativos de la
cuenta GCP del proyecto est\'an disponibles bajo solicitud a trav\'es de
la asociaci\'on.

\section{Desglose detallado de horas-persona}

\begin{table}[H]
    \centering
    \caption{Estimaci\'on de horas dedicadas por miembro y \'area.}
    \label{tab:horas_detalle}
    \begin{tabular}{lcr}
        \toprule
        \textbf{Miembro / Rol} & \textbf{\'Area} & \textbf{Horas estimadas} \\
        \midrule
        Daniel M. (lead)    & Firmware ESP32        & 35 \\
        Daniel M. (lead)    & Servicio RPi          & 15 \\
        Daniel M. (lead)    & Backend FastAPI       & 40 \\
        Daniel M. (lead)    & Frontend              & 20 \\
        Daniel M. (lead)    & SDK + documentaci\'on & 45 \\
        Otros miembros      & ---                   & 0 \\
        \midrule
        \textbf{TOTAL}      &                       & \textbf{155} \\
        \bottomrule
    \end{tabular}
\end{table}

Estas estimaciones son orientativas: la asociaci\'on no implement\'o
sistema formal de \textit{tracking} de horas, por lo que los valores se
apoyan en el an\'alisis de los sellos de tiempo del repositorio
descrito en el cap\'itulo~\ref{cap:sostenibilidad}.

La fila de \textit{otros miembros} en cero no es una omisi\'on, sino un
dato: el repositorio registra \textbf{227 \textit{commits} de un \'unico
autor} a lo largo de los 44 d\'ias naturales de desarrollo activo, sin
ninguna contribuci\'on de terceros al c\'odigo, la documentaci\'on ni el
despliegue. Es la constataci\'on cuantitativa del \textit{bus factor}
igual a uno analizado en la secci\'on~\ref{sec:humanos}.

\section{Fuentes de financiaci\'on recibidas}

\begin{table}[H]
    \centering
    \caption{Fuentes de financiaci\'on recibidas durante el proyecto.}
    \label{tab:financiacion_recibida}
    \begin{tabular}{p{7.5cm}lr}
        \toprule
        \textbf{Fuente} & \textbf{Fecha} & \textbf{Importe} \\
        \midrule
        Universidad de Sevilla --- Proyecto Docente ETSIA
        (expediente 26317, instalaci\'on hidr\'aulica)   & 10/2025 & 900,00\,\euro \\
        Asociaci\'on ESIBot --- material propio y
        aportaciones internas (electr\'onica de control) & 2025--2026 & 366\,\euro \\
        Patrocinios y subvenciones externas              & ---     & 0\,\euro \\
        \midrule
        \textbf{TOTAL FINANCIACI\'ON RECIBIDA}           &         & \textbf{1.266\,\euro} \\
        \bottomrule
    \end{tabular}
\end{table}

Toda la financiaci\'on recibida tiene car\'acter de \textbf{inversi\'on
puntual en material}. Ninguna de las dos fuentes cubri\'o gasto
corriente: los 200\,\euro\ de infraestructura cloud acumulados hasta la
fecha se han sufragado al margen de estas partidas, y no existe ninguna
asignaci\'on comprometida para ejercicios futuros.

\section{Plan de devoluci\'on / cesi\'on de hardware}

\begin{table}[H]
    \centering
    \caption{Destino propuesto del hardware adquirido.}
    \label{tab:plan_hardware}
    \begin{tabular}{p{3.4cm}p{4cm}p{5.2cm}}
        \toprule
        \textbf{Origen financiaci\'on} & \textbf{Material} & \textbf{Destino propuesto} \\
        \midrule
        Proyecto Docente ETSIA (900\,\euro) & Instalaci\'on hidr\'aulica y de potencia (bombas, electrov\'alvulas, v\'alvulas, barriles, fuente DIN, cableado) & Puesta a disposici\'on de la ETSIA como patrimonio universitario. Se propone su permanencia en el almac\'en de la asociaci\'on \'unicamente si la Escuela lo autoriza expresamente. \\
        \midrule
        Asociaci\'on ESIBot (366\,\euro) & Electr\'onica de control, sensores y accesorios & Conservaci\'on en el almac\'en de la asociaci\'on y reaprovechamiento en futuros proyectos. Es material gen\'erico, no espec\'ifico de Demeter. \\
        \midrule
        Aportaciones personales & --- & No procede: no hubo aportaciones personales de material. \\
        \bottomrule
    \end{tabular}
\end{table}

\textbf{Situaci\'on actual}: la totalidad del material, de ambos
or\'igenes, se encuentra depositada en el \textbf{almac\'en de la
asociaci\'on ESIBot}. La instalaci\'on hidr\'aulica no lleg\'o a montarse
en campo, por lo que se conserva sin uso y en su estado original de
recepci\'on.

\textbf{Acci\'on propuesta}: realizar el inventario f\'isico contradictorio
durante la semana 3 del cronograma del Anexo~\ref{anx:plan_entrega},
levantando acta firmada entre la asociaci\'on y la persona designada por
la ETSIA para la recepci\'on, con la tabla~\ref{tab:inventario_hidraulico}
como base de comprobaci\'on unidad a unidad.